% !TeX root = ../paper.tex
\subsection{Data Preparation}
\label{sec:data-preparation}

\subsubsection{Data Extraction and Processing}
% TODO: Describe extraction and processing for text, image, and structured data
% TODO: Detail cleaning, normalization, and quality control procedures


\subsubsection{Pipeline for Each Data Split}

\paragraph{RAG Knowledge Base Pipeline}
\begin{enumerate}
    \item \textbf{Document Chunking:} 
    \begin{itemize}
        \item Strategy: [e.g., sentence-based, paragraph-based, fixed token size]
        \item Chunk size: [e.g., 512 tokens with 50 token overlap]
        \item Preservation: Maintain context and semantic coherence
    \end{itemize}
    
    \item \textbf{Embedding Generation:}
    \begin{itemize}
        \item Model: [e.g., BioBERT, PubMedBERT, sentence-transformers]
        \item Dimension: [e.g., 768-dimensional embeddings]
        \item Storage: Vector database [e.g., FAISS, Pinecone, Chroma]
    \end{itemize}
    
    \item \textbf{Metadata Enrichment:}
    \begin{itemize}
        \item Add disease categories, source information, date, relevance scores
        \item Enable filtering during retrieval
    \end{itemize}
    
    \item \textbf{Index Construction:}
    \begin{itemize}
        \item Build efficient similarity search index
        \item Optimize for low-latency retrieval
    \end{itemize}
\end{enumerate}

\paragraph{Fine-tuning Dataset Pipeline}
\begin{enumerate}
    \item \textbf{Case Selection:}
    \begin{itemize}
        \item Filter for complete cases with diagnosis labels
        \item Ensure disease category diversity and balance
    \end{itemize}
    
    \item \textbf{Input Formatting:}
    \begin{itemize}
        \item Template: \texttt{[Clinical History] + [Symptoms] + [Images] + [Lab Results]}
        \item Multimodal fusion: Combine text and image inputs
        \item Tokenization: Apply model-specific tokenization
    \end{itemize}
    
    \item \textbf{Label Preparation:}
    \begin{itemize}
        \item Format: [e.g., diagnosis name, ICD codes, confidence scores]
        \item Multi-label handling: Support multiple diagnoses per case
    \end{itemize}
    
    \item \textbf{Quality Assurance:}
    \begin{itemize}
        \item Verify label accuracy
        \item Remove ambiguous or incomplete cases
        \item Expert validation: [Specify if clinical expert reviewed data]
    \end{itemize}
\end{enumerate}

\paragraph{Test Dataset Pipeline}
\begin{enumerate}
    \item \textbf{Stratified Sampling:}
    \begin{itemize}
        \item Ensure representation across disease categories
        \item Balance difficulty levels and rare diseases
    \end{itemize}
    
    \item \textbf{Format Consistency:}
    \begin{itemize}
        \item Apply same preprocessing as training data
        \item Maintain multimodal structure
    \end{itemize}
    
    \item \textbf{Ground Truth Verification:}
    \begin{itemize}
        \item Validate diagnosis labels
        \item Document diagnostic rationale for interpretation
    \end{itemize}
\end{enumerate}


\subsubsection{Format Requirements}

\paragraph{RAG Format}
\begin{verbatim}
{
  "chunk_id": "uuid",
  "text": "clinical content...",
  "disease": "disease_name",
  "source": "textbook/article/guideline",
  "embedding": [vector],
  "metadata": {...}
}
\end{verbatim}

\paragraph{Fine-tuning Format}
\begin{verbatim}
{
  "input": {
    "text": "Patient presentation...",
    "image_paths": ["path1.jpg", ...],
    "structured": {demographics, vitals, labs}
  },
  "output": {
    "diagnosis": "disease_name",
    "confidence": 0.95,
    "reasoning": "diagnostic rationale..."
  }
}
\end{verbatim}

\paragraph{Test Format}
Similar to fine-tuning format, with additional ground truth annotations and evaluation metadata.


\subsubsection{Data Partition Summary}

% Add table with actual numbers
\begin{table}[htbp]
\caption{Data Partition Details}
\label{tab:data-partition}
\centering
\begin{tabular}{llll}
\hline
\textbf{Split} & \textbf{Purpose} & \textbf{Size} & \textbf{Ratio} \\
\hline
RAG KB & Retrieval & [X samples] & N/A \\
Train & Fine-tuning & [Y samples] & [e.g., 80\%] \\
Validation & Hyperparameter tuning & [Z samples] & [e.g., 10\%] \\
Test & Evaluation & [W samples] & [e.g., 10\%] \\
\hline
\textbf{Total} & & [Total] & 100\% \\
\hline
\end{tabular}
\end{table}

% Describe your actual partition strategy
\begin{itemize}
    \item \textbf{Partition Strategy:} [e.g., 80-10-10 split for fine-tuning data]
    \item \textbf{Stratification:} Maintained disease category proportions across splits
    \item \textbf{Temporal Split:} [If applicable - e.g., earlier cases for training, recent for testing]
    \item \textbf{No Data Leakage:} Ensured patient-level separation between train/test
\end{itemize}
